\subsection{Dataset} 
In this work, we evaluate our method on two dermatological disease datasets, including ISIC 2019 challenge \cite{combalia2019bcn20000,tschandl2018ham10000} and the Fitzpatrick-17k dataset \cite{groh2021evaluating} . ISIC 2019 challenge contains 25,331 images in 9 diagnostic categories for target labels, and we take gender as our sensitive attribute. The Fitzpatrick-17k dataset contains 16,577 images in 114 skin conditions of target labels and defines skin tone as the sensitive attribute. Next, we apply the data augmentation, including random flipping, rotation, scaling and autoaugment \cite{cubuk2018autoaugment}. After that, we follow the same data split described in \cite{wu2022fairprune} to split the data.
% Finally, we apply the stratified split to the dataset.

% The skin tones are categorized into six levels from 1 to 6. The smaller the value, the lighter the skin. Therefore, we categorize levels 1 to 3 as light skin and 4 to 6 as dark skin

\subsection{Implementation Details and Evaluation Protocol}

% , which determines the contribution of each $CLF$ and classifier to the overall loss function

We employ ResNet18 and VGG-11 as the backbone architectures for our models. The baseline CNN and the multi-exit models are trained for 200 epochs using an SGD optimizer with a batch size of 256 and a learning rate of 1e-2. Each backbone consisted of four internal classifiers ($CLF$s) and a final classifier ($CLF_{f}$). For ResNet18, the internal features are extracted from the end of each residual block, and for VGG-11 the features are extracted from the last four max pooling layers. The loss weight hyperparameter $\alpha$ is selected based on the approach of \cite{kaya2019shallow} and set to $[0.3, 0.45, 0.6, 0.75, 0.9]$ for the multi-exit models and the $\lambda$ is set to 0.01. The confidence threshold $\theta$ of the test set is set to 0.999, based on the best result after performing a grid search on the validation set.

To evaluate the fairness performance of our framework, we adopted the multi-class equalized opportunity (Eopp0 and Eopp1) and equalized odds (Eodd) metrics proposed in \cite{hardt2016equality}. Specifically, we followed the approach of \cite{wu2022fairprune} for calculating these metrics.

\begin{table}[ht]
\centering
\caption{Results of accuracy and fairness of different methods on ISIC 2019 dataset, using gender as the sensitive attribute. The female is the privileged group with higher accuracy by vanilla training. All results are our implementation. }
\setlength{\tabcolsep}{6pt}
\label{Table:compare_isic}
\scriptsize
\begin{tabular}{c c c c c c c c}\toprule 
% \hline
% {\multirow{2}{*}{Method}} & \multicolumn{1}{c@{}}{\multirow{2}{*}{Gender}} &\multicolumn{3}{c@{}}{\text{Accuracy}} & \multicolumn{3}{c}{\text{Fairness}}  \\ 
% \cmidrule(lr){3-5}
% \cmidrule(ll){6-8}
% {} & {} &\text{Precision} & \text{Recall} & \text{F1-score} & \text{Eopp0 $\downarrow$} & \text{Eopp1 $\downarrow$} & \text{Eodd $\downarrow$} \\ 
% \cmidrule(lr){1-1}
% \cmidrule(lr){2-2}
% \cmidrule(lr){3-5}
% \cmidrule(ll){6-8}
\multicolumn{1}{c }{\text{}} & \multicolumn{1}{c }{\text{}} &\multicolumn{3}{c }{\text{Accuracy}} & \multicolumn{3}{c }{\text{Fairness}}  \\ 
\cmidrule(lr){3-5}
\cmidrule(ll){6-8}
\multicolumn{1}{ c }{\text{Method}} & \text{Gender} &\text{Precision} & \text{Recall} & \text{F1-score} & \text{Eopp0 $\downarrow$} & \text{Eopp1 $\downarrow$} & \text{Eodd $\downarrow$} \\ 
\hline
\hline
 {\multirow{4}{*}{ResNet18}} & \text{Female} & 0.793 & 0.721 & 0.746 & \multirow{4}{*}{0.006} & \multirow{4}{*}{0.044} & \multirow{4}{*}{0.022} \\ 
 & \text{Male} & 0.731 & 0.725 & 0.723 &  &  & \\ 
 & \text{Avg. $\uparrow$} & 0.762 & 0.723 & 0.735 &  &  & \\ 
 & \text{Diff. $\downarrow$} & 0.063 & 0.004 & 0.023 &  &  & \\ 
% \hline
% \cmidrule(lr){1-1}
% \cmidrule(lr){2-2}
% \cmidrule(lr){3-5}
% \cmidrule(ll){6-8}

% \multirow{4}{*}{\multicolumn{1}{ c }{ME-ResNet18}} & \text{Female} & 0.748 & 0.724 & 0.733 & \multirow{4}{*}{0.006} & \multirow{4}{*}{0.031} & \multirow{4}{*}{0.016} \\ 
%  & \text{Male} & 0.723 & 0.736 & 0.726 &  &  & \\ 
%  & \text{Avg.} & 0.735 & 0.730 & 0.730 &  &  & \\ 
%  & \text{Diff.} & 0.025 & 0.012 & 0.007 &  &  & \\ 
% \hline
\hline


\multirow{4}{*}{AdvConf} & \text{Female} & 0.755 & 0.738 & 0.741 & \multirow{4}{*}{0.008} & \multirow{4}{*}{0.070} & \multirow{4}{*}{0.037} \\ 
 & \text{Male} & 0.710 & 0.757 & 0.731 &  &  & \\ 
 & \text{Avg. $\uparrow$} & 0.733 & 0.747 & 0.736 &  &  & \\ 
 & \text{Diff. $\downarrow$} & 0.045 & 0.020 & 0.010 &  &  & \\ 
\hline
% \cmidrule(lr){1-1}
% \cmidrule(lr){2-2}
% \cmidrule(lr){3-5}
% \cmidrule(ll){6-8}
\multirow{4}{*}{AdvRev} & \text{Female} & 0.778 & 0.683 & 0.716 & \multirow{4}{*}{0.007} & \multirow{4}{*}{0.059} & \multirow{4}{*}{0.033} \\ 
 & \text{Male} & 0.773 & 0.706 & 0.729 &  &  & \\ 
 & \text{Avg. $\uparrow$} & 0.775 & 0.694 & 0.723 &  &  & \\ 
 & \text{Diff. $\downarrow$} & 0.006 & 0.023 & 0.014 &  &  & \\ 
%  \cmidrule(lr){1-1}
% \cmidrule(lr){2-2}
% \cmidrule(lr){3-5}
% \cmidrule(ll){6-8}
\hline
\multirow{4}{*}{DomainIndep} & \text{Female} & 0.729 & 0.747 & 0.734 & \multirow{4}{*}{0.010} & \multirow{4}{*}{0.086} & \multirow{4}{*}{0.042} \\ 
 & \text{Male} & 0.725 & 0.694 & 0.702 &  &  & \\ 
 & \text{Avg. $\uparrow$} & 0.727 & 0.721 & 0.718 &  &  & \\ 
 & \text{Diff. $\downarrow$} & 0.004 & 0.053 & 0.031 &  &  & \\ 
% \hline
\hline
\multirow{4}{*}{HSIC} & \text{Female} & 0.744 & 0.660 & 0.696 & \multirow{4}{*}{0.008} & \multirow{4}{*}{0.042} & \multirow{4}{*}{0.020} \\ 
 & \text{Male} & 0.718 & 0.697 & 0.705 &  &  & \\ 
 & \text{Avg. $\uparrow$} & 0.731 & 0.679 & 0.700 &  &  & \\ 
 & \text{Diff. $\downarrow$} & 0.026 & 0.037 & 0.009 &  &  & \\ 
% \cmidrule(lr){1-1}
% \cmidrule(lr){2-2}
% \cmidrule(lr){3-5}
% \cmidrule(ll){6-8}
% % \hline

% \multirow{4}{*}{\multicolumn{1}{ c }{ME-HSIC}} & \text{Female} & 0.733 & 0.707 & 0.716 & \multirow{4}{*}{0.007} & \multirow{4}{*}{0.034} & \multirow{4}{*}{0.018} \\ 
%  & \text{Male} & 0.713 & 0.707 & 0.707 &  &  & \\ 
%  & \text{Avg.} & 0.723 & 0.707 & 0.712 &  &  & \\ 
%  & \text{Diff.} & 0.020 & 0.000 & 0.009 &  &  & \\ 
% \hline
\hline
% \multicolumn{1}{ c }{HSIC} & 0.74 & 15 & 630 & 0.70 & 15 & 618 & -\\ 
\multicolumn{1}{ c }{\multirow{4}{*}{MFD}} & \text{Female} & 0.770 & 0.697 & 0.726 & \multirow{4}{*}{\textbf{0.005}} & \multirow{4}{*}{0.051} & \multirow{4}{*}{0.024} \\ 
 & \text{Male} & 0.772 & 0.726 & 0.744 &  &  & \\ 
 & \text{Avg. $\uparrow$} & 0.771 & 0.712 & 0.735 &  &  & \\ 
 & \text{Diff. $\downarrow$} & 0.002 & 0.029 & 0.018 &  &  & \\ 
% \hline
% \cmidrule(lr){1-1}
% \cmidrule(lr){2-2}
% \cmidrule(lr){3-5}
% \cmidrule(ll){6-8}

% \multirow{4}{*}{\multicolumn{1}{ c }{ME-MFD}} & \text{Female} & 0.733 & 0.698 & 0.711 & \multirow{4}{*}{0.006} & \multirow{4}{*}{0.024} & \multirow{4}{*}{0.012} \\ 
%  & \text{Male} & 0.772 & 0.713 & 0.739 &  &  & \\ 
%  & \text{Avg.} & 0.752 & 0.706 & 0.725 &  &  & \\ 
%  & \text{Diff.} & 0.039 & 0.015 & 0.028 &  &  & \\ 
\hline
\hline
% \multicolumn{1}{ c }{MFD} & 0.06 & 7 & 21 & 0.06 & 5 & 27 & -\\ 
% \multirow{4}{*}{\multicolumn{1}{ c }{DomainIndep}} & \text{Female} & 0.729 & 0.747 & 0.734 & \multirow{4}{*}{0.010} & \multirow{4}{*}{0.086} & \multirow{4}{*}{0.042} \\ 
%  & \text{Male} & 0.725 & 0.694 & 0.702 &  &  & \\ 
%  & \text{Avg.} & 0.727 & 0.721 & 0.718 &  &  & \\ 
%  & \text{Diff.} & 0.004 & 0.053 & 0.031 &  &  & \\ 
% \hline
% \multicolumn{1}{ c }{DomainIndep} & 0.06 & 7 & 21 & 0.06 & 5 & 27 & -\\ 
{\multirow{4}{*}{FairPrune}} & \text{Female} & 0.776 & 0.711 & 0.734 & \multirow{4}{*}{0.007} & \multirow{4}{*}{0.026} & \multirow{4}{*}{0.014} \\ 
 & \text{Male} & 0.721 & 0.725 & 0.720 &  &  & \\ 
 & \text{Avg. $\uparrow$} & 0.748 & 0.718 & 0.727 &  &  & \\ 
 & \text{Diff. $\downarrow$} & 0.055 & 0.014 & 0.014 &  &  & \\ 
\hline
% \cmidrule(lr){1-1}
% \cmidrule(lr){2-2}
% \cmidrule(lr){3-5}
% \cmidrule(ll){6-8}
% \multicolumn{1}{ c }{FairPrune} & 0.06 & 7 & 21 & 0.06 & 5 & 27 & -\\ 
\multicolumn{1}{ c }{\multirow{4}{*}{ME-FairPrune}} & \text{Female} & 0.770 & 0.723 & 0.742 & \multirow{4}{*}{0.006} & \multirow{4}{*}{\textbf{0.020}} & \multirow{4}{*}{\textbf{0.010}} \\ 
 & \text{Male} & 0.739 & 0.728 & 0.730 &  &  & \\ 
 & \text{Avg. $\uparrow$} & 0.755 & 0.725 & 0.736 &  &  & \\ 
 & \text{Diff. $\downarrow$} & 0.032 & 0.005 & 0.012 &  &  & \\ 
 % \hline
\bottomrule
\end{tabular}
\end{table}


% ver 1
% \begin{table}[ht]
% \centering
% \label{Table:compare_isic}
% \scriptsize
% \begin{tabular}{N N N N N N N N}\toprule
% \multicolumn{1}{N }{\text{}} & \multicolumn{1}{c }{\text{}} &\multicolumn{3}{c }{\text{Accuracy}} & \multicolumn{3}{c }{\text{Fairness}}  \\ 
% \cmidrule(lr){3-5}
% \cmidrule(ll){6-8}
% \multicolumn{1}{ c }{\text{Method}} & \text{Gender} &\text{Precision} & \text{Recall} & \text{F1-score} & \text{Eopp0} & \text{Eopp1} & \text{Eodd} \\ 
% % \cmidrule(lr){1-1}
% % \cmidrule(lr){2-2}
% % \cmidrule(lr){3-5}
% % \cmidrule(ll){6-8}
% \hline
% \multirow{4}{*}{\multicolumn{1}{ c }{ResNet18}} & \text{Female} & 0.793 & 0.721 & 0.746 & \multirow{4}{*}{0.006} & \multirow{4}{*}{0.044} & \multirow{4}{*}{0.022} \\ 
%  & \text{Male} & 0.731 & 0.725 & 0.723 &  &  & \\ 
%  & \text{Avg.} & 0.762 & 0.723 & 0.735 &  &  & \\ 
%  & \text{Diff.} & 0.063 & 0.004 & 0.023 &  &  & \\ 
% \hline

% \multirow{4}{*}{\multicolumn{1}{ c }{AdvConf}} & \text{Female} & 0.755 & 0.738 & 0.741 & \multirow{4}{*}{0.008} & \multirow{4}{*}{0.070} & \multirow{4}{*}{0.037} \\ 
%  & \text{Male} & 0.710 & 0.757 & 0.731 &  &  & \\ 
%  & \text{Avg.} & 0.733 & 0.747 & 0.736 &  &  & \\ 
%  & \text{Diff.} & 0.045 & 0.020 & 0.010 &  &  & \\ 
% % \hline
% \cmidrule(lr){1-1}
% \cmidrule(lr){2-2}
% \cmidrule(lr){3-5}
% \cmidrule(ll){6-8}
% \multirow{4}{*}{\multicolumn{1}{ c }{AdvRev}} & \text{Female} & 0.778 & 0.683 & 0.716 & \multirow{4}{*}{0.007} & \multirow{4}{*}{0.059} & \multirow{4}{*}{0.033} \\ 
%  & \text{Male} & 0.773 & 0.706 & 0.729 &  &  & \\ 
%  & \text{Avg.} & 0.775 & 0.694 & 0.723 &  &  & \\ 
%  & \text{Diff.} & 0.006 & 0.023 & 0.014 &  &  & \\ 
% % \hline
% \multirow{4}{*}{\multicolumn{1}{ c }{DomainIndep}} & \text{Female} & 0.729 & 0.747 & 0.734 & \multirow{4}{*}{0.010} & \multirow{4}{*}{0.086} & \multirow{4}{*}{0.042} \\ 
%  & \text{Male} & 0.725 & 0.694 & 0.702 &  &  & \\ 
%  & \text{Avg.} & 0.727 & 0.721 & 0.718 &  &  & \\ 
%  & \text{Diff.} & 0.004 & 0.053 & 0.031 &  &  & \\ 
% \hline

% \multirow{4}{*}{\multicolumn{1}{ c }{HSIC}} & \text{Female} & 0.744 & 0.660 & 0.696 & \multirow{4}{*}{0.008} & \multirow{4}{*}{0.042} & \multirow{4}{*}{0.019} \\ 
%  & \text{Male} & 0.718 & 0.697 & 0.705 &  &  & \\ 
%  & \text{Avg.} & 0.731 & 0.679 & 0.700 &  &  & \\ 
%  & \text{Diff.} & 0.026 & 0.037 & 0.009 &  &  & \\ 
% \hline

% % \multicolumn{1}{ c }{HSIC} & 0.74 & 15 & 630 & 0.70 & 15 & 618 & -\\ 
% \multirow{4}{*}{\multicolumn{1}{ c }{MFD}} & \text{Female} & 0.770 & 0.697 & 0.727 & \multirow{4}{*}{0.005} & \multirow{4}{*}{0.051} & \multirow{4}{*}{0.024} \\ 
%  & \text{Male} & 0.772 & 0.726 & 0.744 &  &  & \\ 
%  & \text{Avg.} & 0.771 & 0.712 & 0.736 &  &  & \\ 
%  & \text{Diff.} & 0.002 & 0.029 & 0.017 &  &  & \\ 
% \hline

% % \multicolumn{1}{ c }{MFD} & 0.06 & 7 & 21 & 0.06 & 5 & 27 & -\\ 
% % \multirow{4}{*}{\multicolumn{1}{ c }{DomainIndep}} & \text{Female} & 0.729 & 0.747 & 0.734 & \multirow{4}{*}{0.010} & \multirow{4}{*}{0.086} & \multirow{4}{*}{0.042} \\ 
% %  & \text{Male} & 0.725 & 0.694 & 0.702 &  &  & \\ 
% %  & \text{Avg.} & 0.727 & 0.721 & 0.718 &  &  & \\ 
% %  & \text{Diff.} & 0.004 & 0.053 & 0.031 &  &  & \\ 
% % \hline
% % \multicolumn{1}{ c }{DomainIndep} & 0.06 & 7 & 21 & 0.06 & 5 & 27 & -\\ 
% \multirow{4}{*}{\multicolumn{1}{ c }{FairPrune}} & \text{Female} & 0.776 & 0.711 & 0.734 & \multirow{4}{*}{0.007} & \multirow{4}{*}{0.026} & \multirow{4}{*}{0.014} \\ 
%  & \text{Male} & 0.721 & 0.725 & 0.720 &  &  & \\ 
%  & \text{Avg.} & 0.748 & 0.718 & 0.727 &  &  & \\ 
%  & \text{Diff.} & 0.055 & 0.014 & 0.014 &  &  & \\ 
% \hline
% % \multicolumn{1}{ c }{FairPrune} & 0.06 & 7 & 21 & 0.06 & 5 & 27 & -\\ 
% \multirow{4}{*}{\multicolumn{1}{ c }{ME-FairPrune}} & \text{Female} & 0.770 & \textbf{0.723} & 0.742 & \multirow{4}{*}{0.006} & \multirow{4}{*}{\textbf{0.020}} & \multirow{4}{*}{\textbf{0.010}} \\ 
%  & \text{Male} & 0.739 & 0.728 & 0.730 &  &  & \\ 
%  & \text{Avg.} & 0.755 & 0.725 & 0.736 &  &  & \\ 
%  & \text{Diff.} & 0.032 & 0.005 & 0.012 &  &  & \\ 
% \bottomrule
% \end{tabular}
% \end{table}


% \begin{tabular}{SSSSSSSS} 
% \toprule
% \multicolumn{7}{c}{Accuracy} {Fairness}\\
% \cmidrule(r){3-5} \cmidrule(lr){6-8}
%     {Method} & {Skin Tone} & {Precision} & {Recall} & {F1-score} & {Eopp0} & {Eopp1} & {Eodd} \\ \midrule
%     1  & 16.128 & +8.872 & 16.128 & 1.402 & 1.373 & -146.6 & -137.6 \\
%     2  & 3.442  & -2.509 & 3.442  & 0.299 & 0.343 & 133.2  & 152.4  \\
%     3  & 1.826  & -0.363 & 1.826  & 0.159 & 0.119 & 168.5  & -161.1 \\
%     4  & 0.993  & -0.429 & 0.993  & 0.086 & 0.08  & 25.6   & 90     \\ \midrule
%     5  & 1.29   & +0.099 & 1.29   & 0.112 & 0.097 & -175.6 & -114.7 \\
%     6  & 0.483  & -0.183 & 0.483  & 0.042 & 0.063 & 22.3   & 122.5  \\
%     7  & 0.766  & -0.475 & 0.766  & 0.067 & 0.039 & 141.6  & -122   \\
%     8  & 0.624  & +0.365 & 0.624  & 0.054 & 0.04  & -35.7  & 90     \\ \midrule
%     9  & 0.641  & -0.466 & 0.641  & 0.056 & 0.045 & 133.3  & -106.3 \\
%     10 & 0.45   & +0.421 & 0.45   & 0.039 & 0.034 & -69.4  & 110.9  \\
%     11 & 0.598  & -0.597 & 0.598  & 0.052 & 0.025 & 92.3   & -109.3 \\ \bottomrule
% \end{tabular}


% The Eopp0, Eopp1 and the Eodd are computed as below equations, where $TPR_{k}^{a}$, $TNR_{k}^{a}$ and $FPR_{k}^{a}$ are the True Positive Rate, True Negative Rate and the False Positive Rate respectively of target class $k$ and sensitive attribute $a$: 
% \begin{equation}
% \label{eq.EOpp0_equation}
%     EOpp0 = \sum_{k=1}^{K} |TNR_{k}^{1} - TNR_{k}^{0}|, \: EOpp1 = \sum_{k=1}^{K} |TPR_{k}^{1} - TPR_{k}^{0}|
% \end{equation}

% % \begin{equation}
% % \label{eq.EOpp1_equation}
% %     EOpp1 = \sum_{k=1}^{K} |TPR_{k}^{1} - TPR_{k}^{0}|
% % \end{equation}

% \begin{equation}
% \label{eq.EO_equation}
%     EOdd = \sum_{k=1}^{K} |TPR_{k}^{1} - TPR_{k}^{0} + FPR_{k}^{1} - FPR_{k}^{0}|
% \end{equation}

% \subsection{Baselines}
% We compared our models with several baselines. \emph{Vanilla} is a standard CNN without applying a debiased strategy. \emph{AdvConf} \cite{alvi2018turning} and \emph{AdvRev} \cite{zhang2018mitigating} are adversarial training based de-bias method. \emph{HSIC} \cite{quadrianto2019discovering} and \emph{MFD} \cite{jung2021fair} are both regularization-based de-bias method which add the regularization term at loss function to add the fairness constraint in the training process. \emph{DomainIndep} \cite{wang2020towards} uses multiple classifier to explicitly separate information of each group. \emph{FairPrune} \cite{wu2022fairprune} achieves fair dermatology disease classification by pruning the neurons with unfairly favorable to certain group. 

% \subsection{Evaluation Metrics}
% We evaluate our model by multi-class equalized opportunity (Eopp) and equalized odds (Eodd) to quantify the fairness performance \cite{hardt2016equality}. The equalized opportunity has two metrics to evaluate the correct prediction. Eopp0 indicates the difference between the two groups' true negative rate (TNR), and Eopp1 illustrates the difference between the two groups' true positive rate (TPR). Equalized odds is the summation of two groups' true positive rate difference and false positive rate difference, which can also address unfair wrong prediction of the model. The Eopp and the Eodd are computed as below equations, where $TPR_{k}^{a}$, $TNR_{k}^{a}$ and $FPR_{k}^{a}$ are the True Positive Rate, True Negative Rate and the False Positive Rate respectively of target class $k$ and sensitive attribute $a$: 
% \begin{equation}
% \label{eq.EOpp0_equation}
%     EOpp0 = \sum_{k=1}^{K} |TNR_{k}^{1} - TNR_{k}^{0}|, \: EOpp1 = \sum_{k=1}^{K} |TPR_{k}^{1} - TPR_{k}^{0}|
% \end{equation}

% % \begin{equation}
% % \label{eq.EOpp1_equation}
% %     EOpp1 = \sum_{k=1}^{K} |TPR_{k}^{1} - TPR_{k}^{0}|
% % \end{equation}

% \begin{equation}
% \label{eq.EO_equation}
%     Eodd = \sum_{k=1}^{K} |TPR_{k}^{1} - TPR_{k}^{0} + FPR_{k}^{1} - FPR_{k}^{0}|
% \end{equation}



% \subsection{Dataset} 
% In this work, we evaluated our method on two dermatology disease classification datasets, including ISIC 2019 challenge \cite{combalia2019bcn20000,tschandl2018ham10000} and the Fitzpatrick-17k dataset \cite{groh2021evaluating} . ISIC 2019 challenge contains 25,331 images in 9 diagnostic categories. We took the gender labels as our sensitive attribute.
% The Fitzpatrick-17k dataset contains 16,577 images in 114 skin conditions of target labels and defines the skin tone labels as the sensitive attributes. The skin tones are categorized into six levels from 1 to 6. The smaller the value, the lighter the skin. Therefore, we categorize levels 1 to 3 as light skin and 4 to 6 as dark skin. In this work, we applied the data augmentation technique and follow the dataset split described in \cite{wu2022fairprune}.

% % We used ResNet18 and VGG11 as the backbone, and applied the data augmentation technique and follow the dataset split described in \cite{wu2022fairprune}. The CNN baseline and multi-exit models were both trained for 200 epochs using an SGD optimizer with a learning rate of 1e-2 and a batch size of 256. The hyperparameter $\alpha$, which is the loss weight of each $IC$, was determined using the setting from \cite{kaya2019shallow} as $[0.3, 0.45, 0.6, 0.75, 0.9]$ for the multi-exit models. The early exit confidence threshold for test set was set to 0.999 after performing a grid search on the validation set.

% \subsection{Baselines}
% We compared our models with several baselines. \emph{Vanilla} is a standard CNN without applying a debiased strategy. \emph{AdvConf} \cite{alvi2018turning} and \emph{AdvRev} \cite{zhang2018mitigating} are adversarial training based de-bias method. \emph{HSIC} \cite{quadrianto2019discovering} and \emph{MFD} \cite{jung2021fair} are both regularization-based de-bias method which add the regularization term at loss function to add the fairness constraint in the training process. \emph{DomainIndep} \cite{wang2020towards} uses multiple classifier to explicitly separate information of each group. \emph{FairPrune} \cite{wu2022fairprune} achieves fair dermatology disease classification by pruning the neurons with unfairly favorable to certain group. 

% \subsection{Evaluation Metrics}
% We evaluate our model by multi-class equalized opportunity (EOpp) and equalized odds (EOdd) to quantify the fairness performance \cite{hardt2016equality}. The equalized opportunity has two metrics to evaluate the correct prediction. EOpp0 indicates the difference between the two groups' true negative rate (TNR), and EOpp1 illustrates the difference between the two groups' true positive rate (TPR). Equalized odds is the summation of two groups' true positive rate difference and false positive rate difference, which can also address unfair wrong prediction of the model. The EOpp and the EOdd are computed as below equations, where $TPR_{k}^{a}$, $TNR_{k}^{a}$ and $FPR_{k}^{a}$ are the True Positive Rate, True Negative Rate and the False Positive Rate respectively of target class $k$ and sensitive attribute $a$: 
% \begin{equation}
% \label{eq.EOpp0_equation}
%     EOpp0 = \sum_{k=1}^{K} |TNR_{k}^{1} - TNR_{k}^{0}|, \: EOpp1 = \sum_{k=1}^{K} |TPR_{k}^{1} - TPR_{k}^{0}|
% \end{equation}

% % \begin{equation}
% % \label{eq.EOpp1_equation}
% %     EOpp1 = \sum_{k=1}^{K} |TPR_{k}^{1} - TPR_{k}^{0}|
% % \end{equation}

% \begin{equation}
% \label{eq.EO_equation}
%     Eodd = \sum_{k=1}^{K} |TPR_{k}^{1} - TPR_{k}^{0} + FPR_{k}^{1} - FPR_{k}^{0}|
% \end{equation}
